% 示例：基于 MedViT 的布匹织物分类方法（LaTeX 文档框架）
\documentclass[UTF8]{ctexart}
\usepackage{amsmath,amssymb}
\usepackage{graphicx}
\usepackage{booktabs}
\usepackage{hyperref}

\title{基于 MedViT 混合架构的布匹织物分类方法及应用改进}
\author{示例作者}
\date{\today}

\begin{document}
\maketitle

\section{技术方案}

\subsection{预处理}
对织物图像进行尺寸标准化（$224\times 224$ 像素），采用 ImageNet 均值与标准差归一化。

\subsection{MedViT 混合架构}
网络包含局部特征感知块（LFP）与全局特征感知块（GFP），其中：
\begin{itemize}
  \item LFP：采用扩张邻域注意力（DiNA），复杂度 $O(N\times k)$；
  \item GFP：采用 E-MHSA 与 KAN 层，替代传统 MLP。
\end{itemize}

\subsection{损失与训练}
采用交叉熵损失：
\[
  L = -\frac{1}{N}\sum_{i=1}^{N}\sum_{c=1}^{C} y_{i,c}\log(p_{i,c})
\]
优化器为 AdamW，学习率 $1\times 10^{-4}$，权重衰减 $0.05$，余弦退火 300 轮。

\section{实验结果（示例）}

\begin{table}[htbp]
  \centering
  \caption{TextileNet-fabric 织物分类 Top-1 准确率（\%）}
  \begin{tabular}{lcc}
    \toprule
    方法        & Top-1 & 参数量 (M) \\
    \midrule
    ResNet50    & 65.28 & 25.6       \\
    MedViTV2-S  & 72.1  & 28.4       \\
    \bottomrule
  \end{tabular}
\end{table}

\end{document}
